\section{Discussion, limitations, and declarations}\label{sec:discussion}

\subsection{What exact JC data recover}

The classification gives a concrete answer to the biological identifiability question within the stated class.  An exact generic site-pattern distribution determines the homeomorphism-reduced arrangement of reticulated components, whether each local component is a cycle or a theta blob, the placement of every labelled descendant block, and its position in the ordered port words.  Ordinary triangle redirection is the only topology change that can remain observationally compatible on a full-dimensional model region.  Thus the presence and location of a close reticulation can be identifiable even when its direction inside a three-cycle is not generically determined modulo that move.

This conclusion is stronger than algebraic distinguishability alone.  The theorem classifies one-sided full-dimensional stochastic containment on the open biological parameter domain.  It therefore excludes both equal-dimensional non-$\Tmove$ overlap and a lower-dimensional topology whose generic distributions could all be explained by a larger competing topology.  Conversely, the three ordinary triangle orientations are shown to share a regular open germ; equality of their complete stochastic images is neither asserted nor needed.  In particular, the structural reconstruction returns a canonical topology modulo $\Tmove$, not the input-specific list of every redirected orientation whose stochastic image contains the observed distribution.

The sharpness result explains why the strong convention is substantive.  A topology in $\TCw$ may admit a biologically appealing tree-child rooting while also admitting rootings that expose a vertex with two reticulation children.  The Theta pair shows that this extra flexibility permits an exchange of labelled descendant components to become observationally invisible on a full-dimensional set.  Requiring every admissible rooting to be tree-child removes precisely this mechanism in the theorem's class.

\subsection{Why no triangle-count hypothesis remains}

The final theorem covers all binary standard semi-directed $\TCs$ level-2 networks.  No additional hypothesis on the number of triangles is hidden in the finite atlas.  The reason is structural: a level-2 blob with two triangles must be the $(1,2,2)$ theta, or $K_4$ minus an edge, and \cref{thm:automatic-triangle-bound} shows that this graph has no tree-child rooting.  In fact, the exclusion already holds in the larger class $\TCw$.  The four oriented theta cores and the cycle core therefore exhaust every local factor that can occur in the positive theorem.

The theorem is specific to the displayed-tree JC model with uniform root distribution, strictly positive finite branch lengths, and inheritance probabilities strictly between zero and one.  Boundary parameters can collapse edges, remove reticulation choices, or create containments absent from the open model.  Models with coalescent processes, site-rate variation, nonstationary substitution, or richer group-based edge parameters require separate analysis.  Level greater than two also permits generator cores and triangle interactions not covered by the present structural lemma.

\subsection{Statistical and computational limitations}

Generic identifiability is an infinite-data property.  It does not imply that all identifiable features can be estimated accurately from finite sequences, especially near the algebraic exceptional locus or near a triangle model intersection.  The semialgebraic work on three-leaf triangles shows that substantial full-dimensional overlap may coexist with regions where the triangle direction is distinguishable.  Developing statistically stable estimators for the recoverable canonical class is an important next step.

The reconstruction theorem is constructive but not offered as an efficient implementation.  Its optional enumeration is the structural $\Tmove$-equivalence class; deciding which of those orientations contain one particular distribution is a separate semialgebraic membership problem.  Testing all leaf splits directly is exponential, and exact algebraic arithmetic is expensive.  The finite atlas suggests practical invariant screens, but a statistically and computationally optimized inference method lies outside the present scope.

\subsection{Open problems}

Natural questions remaining within phylogenetic algebraic statistics include the following.  Does an analogous canonical classification hold at level greater than two?  Can the generic bridge and local-tensor reconstruction be made polynomial-time under an oracle model for exact probabilities?  Which parts of the theorem survive under rate variation or nonstationary models?  Finally, can the global bridge/local-tensor extraction be combined with the known semialgebraic triangle inequalities to determine, for one exact distribution, precisely which redirected orientations are pointwise compatible and to give a statistically useful confidence region for the canonical topology class?

\subsection*{Acknowledgements}

The author thanks the authors of the cited phylogenetic-network identifiability papers for making their preprints and computational materials publicly available.  The author also thanks the adversarial AI-assisted review processes and separately implemented automated checks that exposed and corrected several earlier incomplete arguments before this version was prepared.  No institutional funding supported this work.

\subsection*{Use of artificial intelligence}

OpenAI ChatGPT, including GPT-5.6 Pro, and Anthropic Claude were used during exploration and revision.  These systems generated candidate arguments, symbolic-computation code, manuscript drafts, and adversarial review prompts.  Distinct implementations were used to rederive the load-bearing finite calculations, and failed or incomplete intermediate claims were retained in the archived research logs.  Artificial-intelligence systems are not authors.  The human author personally reviewed the final mathematical statements, citations, manuscript, and release artifacts and accepts full responsibility for their accuracy.  Prompts, logs, source code, exact certificates, and verification transcripts are archived with the reproducibility package.

\subsection*{Author contributions}

Alec Kriebel conceived the project, selected and refined the research question, directed the computational and proof audits, checked the final arguments and citations, prepared the release, and accepts responsibility for the work in its entirety.

\subsection*{Competing interests}

The author declares no competing interests.

\subsection*{Data and code availability}

No empirical biological data were used.  Complete LaTeX and TikZ source, primitive graph generators, exact Fourier and polynomial verifiers, finite certificates, independent replay programs, environment specifications, manifests, and successful clean verification transcripts accompany this manuscript.  The versioned supplementary archive accompanying this submission is self-contained and includes its own SHA-256 manifest and verification transcripts.  A persistent repository identifier should be inserted in the journal metadata and this statement after the author deposits the final, checksum-matched archive.  No persistent identifier is claimed in the present local release candidate.
